\begin{explanationbox}
\textbf{\textcolor{CExplanation}{一句话直觉}}
三角形三个内角的余弦值之间并非独立，它们的平方和与乘积通过一个简洁的常数$1$联系。

\medskip
\textbf{\textcolor{CExplanation}{核心拆解}}
\begin{description}[style=nextline, leftmargin=2.8em, labelsep=0.5em, itemsep=0.45em]
    \item[\textbf{要点一（对称结构）：}] 恒等式关于$A,B,C$轮换对称，体现了三角形内角的对等地位。
    \item[\textbf{要点二（乘积修正）：}] 平方和并非恒为$1$，而是用$2\cos A\cos B\cos C$作为修正项，当某角为直角时该项消失。
    \item[\textbf{要点三（内角约束）：}] 导出的关键是使用条件$A+B+C=\pi$，将三角函数的和差转化为单角函数。
\end{description}

\medskip
\textbf{\textcolor{CExplanation}{几何本质（恒等变换）}}
该恒等式并非直接揭示某个几何图形的度量性质，而是源于代数化的三角恒等变形。它刻画了“三内角余弦平方和”这一对称量在三角形约束下的必然取值。

\medskip
\textbf{\textcolor{CExplanation}{代数意义（降次与化简）}}
左边是余弦平方和，可通过降次公式转化为倍角余弦和，右边乘积也可视为三角方程组中的纽带。它常用于将高次三角式降为二次，或将乘积转化为平方和，简化计算。

\medskip
\textbf{\textcolor{CExplanation}{考点价值（高频应用）}}
\begin{itemize}[leftmargin=*, itemsep=0.5em, topsep=0.4em]
    \item \textbf{考法一（直接求值）：} 已知两个角的余弦或乘积关系，利用恒等式快速求出平方和或缺少的余弦值。
    \item \textbf{考法二（恒等证明）：} 在证明其他三角恒等式时，将该式作为已知结论进行代换，减少推导步骤。
\end{itemize}

\medskip
\textbf{\textcolor{CExplanation}{顿悟点}}
\textbf{对称与修正：}记住$A=B=C=\frac{\pi}{3}$时，各项均为$\frac14$，平方和为$\frac34$，右边$1-2\cdot\frac18=\frac34$。这个特例能帮助检查符号和系数。

\medskip
\textbf{\textcolor{CExplanation}{使用场景}}
\begin{itemize}[leftmargin=*, itemsep=0.5em, topsep=0.4em]
    \item \textbf{场景一（已知乘积求平方和）：} 题目给出$\cos A\cos B\cos C$的值，直接代入即得余弦平方和。
    \item \textbf{场景二（已知两角余弦求第三角相关式）：} 结合内角关系，可借助恒等式确定第三角的余弦或平方，避免复杂求角运算。
\end{itemize}
\end{explanationbox}
